\chapter{Ablations: What Each Piece Is Worth}
\label{ch:ablations}

\Cref{ch:trajectory} showed the architecture with every piece switched on. This
chapter asks the harder question: what is each piece actually worth? The honest answer
is built from the contrasts the run already supplies, one component-level switch that
was genuinely toggled, and a clear ledger of which lesions remain to be run. We do not
pretend to a brute-force sweep that was not performed, and the reason that is
defensible is the subject of the first section.

\begin{intuition}[Ablation by measurement, not by brute force]
The textbook way to price a component is to remove it, re-run everything, and measure
the drop. That is expensive: a system with a dozen parts would need a dozen full
re-runs. \caem{} takes a different route wherever it can. Many components leave a
\emph{direct trace} in the trajectory already measured: the value of the router is the
per-tier accuracy and tier shares of \Cref{ch:trajectory}; the value of
self-improvement is the cold-start-versus-trained contrast; the value of the forgetting
guard is what it did when it fired. Where a component's worth is already visible in a
number we have, re-running a lesion to re-measure it would be redundant. Where it is
not, a lesion is warranted. This chapter prices each piece by the cheapest sound
evidence available, and is explicit about the one lesion whose clean counterfactual is
still owed.
\end{intuition}

\section{The ablation registry, and the claim it defends}
\label{sec:abl-registry}

\caem{} ships a small, deliberate ablation registry rather than a large sweep. The
registry pairs each lesion with the specific thesis claim it is meant to defend, and
prunes any lesion whose claim is already settled by a direct metric.

\begin{precise}[Four entries, each tied to a claim]
The registry holds a full-system reference and three lesions. The reference, named
\emph{full}, applies no mutation and is the anchor every delta is measured against. The
three lesions are: \emph{no-self-improvement}, which runs the full pipeline with no
fine-tuning so the system stays at its cold-start weights; \emph{no-forgetting-guard},
which sets the retention tolerance to zero so every update commits regardless of how
much it forgets; and \emph{no-retroverify}, which skips the per-cycle re-scoring of
stored memory. Each defends a numbered claim: no-self-improvement and no-retroverify
defend the self-improvement claim, no-retroverify additionally defends the memory-purity
claim, and no-forgetting-guard defends the no-catastrophic-forgetting claim. The
memory-routing claim is defended without a lesion at all, by the per-tier shares and
accuracies already in \Cref{ch:trajectory}.
\end{precise}

\begin{figure}[t]
\centering
\begin{tikzpicture}[
  font=\footnotesize\sffamily,
  claim/.style={rectangle, rounded corners=4pt, draw=cbFormula!70!black, fill=cbFormula!10,
    line width=1pt, align=center, inner sep=4pt, text width=3.0cm, minimum height=0.9cm},
  metric/.style={rectangle, rounded corners=3pt, draw=cbExample!70!black, fill=cbExample!12,
    line width=0.9pt, align=center, inner sep=3pt, text width=3.6cm, minimum height=0.8cm},
  lesion/.style={rectangle, rounded corners=3pt, draw=cbPitfall!65!black, fill=cbPitfall!10,
    line width=0.9pt, align=center, inner sep=3pt, text width=3.6cm, minimum height=0.8cm},
  flow/.style={-{Stealth[length=2.2mm]}, line width=0.9pt, draw=cbInk!70},
]
  \node[claim] (c1) at (0,3.0)  {memory routing};
  \node[claim] (c2) at (0,1.5)  {memory purity};
  \node[claim] (c3) at (0,0.0)  {self-improvement};
  \node[claim] (c4) at (0,-1.5) {no catastrophic\\forgetting};
  \node[metric] (m1) at (5.6,3.0) {per-tier shares + accuracy\\(\Cref{ch:trajectory}, measured)};
  \node[metric] (m2) at (5.6,1.5) {purity validation\\(\Cref{ch:theory}, measured)};
  \node[metric] (m3) at (5.6,0.4) {cold-start vs trained\\(\Cref{ch:benchmarks}, measured)};
  \node[lesion] (l1) at (5.6,-0.7) {no-retroverify\\(lesion, deferred)};
  \node[lesion] (l2) at (5.6,-1.7) {no-forgetting-guard\\(observed at C4)};
  \draw[flow] (c1) -- (m1);
  \draw[flow] (c2) -- (m2);
  \draw[flow] (c2) -- (l1);
  \draw[flow] (c3) -- (m3);
  \draw[flow] (c3) -- (l1);
  \draw[flow] (c4) -- (l2);
\end{tikzpicture}
\caption[How each claim is defended]{The ablation registry maps each thesis claim to
its cheapest sound evidence. Three claims are settled by a metric already in hand
(green): routing by the per-tier trajectory, purity by the validation measurement,
self-improvement by the cold-start-versus-trained contrast. Two lean on a lesion (red):
no-forgetting-guard, observed naturally when the guard fired at the fourth cycle, and
no-retroverify, whose clean counterfactual is registered but not yet run.
(Schematic.)}
\label{fig:abl-map}
\end{figure}

\begin{bythenumbers}[Why the sweep is small]
An earlier design listed more than twenty candidate lesions. They were pruned against
two questions: is the claim already defended by a direct metric, and is the evidence
redundant with a lesion already kept? A lesion that removed the storage gate was
dropped because the purity validation measures the gate's effect directly; a lesion
that removed the cheap memory tier was dropped because the per-tier share trajectory
already shows its usage; lesions on individual verifier signals were dropped because
the signal correlation matrix shows each carries non-redundant information. What
survived is the minimal set whose claims no existing number answers. This is a
deliberate choice to spend evidence where it is missing, not to manufacture re-runs
where it already exists.
\end{bythenumbers}

\section{Ablating the whole architecture, and the self-improvement slice}
\label{sec:abl-architecture}

The largest ablation is the one already drawn in \Cref{ch:benchmarks}: the eight-baseline
panel is the architecture removed in pieces, and the decomposition is self-improvement
removed on its own.

\begin{precise}[The panel is a component-removal ladder]
Read as ablations, the baselines of \Cref{ch:benchmarks} price the architecture from the
bottom up. The zero-shot floor is everything removed: no routing, no retrieval, no
memory, no verifier, no fine-tuning. Unconditional retrieval adds back grounding and
nothing else. Vanilla fine-tuning adds back training-time learning and nothing else.
The gap from the floor to full \caem{} is therefore the joint worth of routing,
verification, memory, and self-improvement; and the gap from vanilla fine-tuning to
full \caem{} isolates the deployment-time machinery from the training-time machinery. On
the pooled panel the floor-to-cold-start step, which adds the three tiers, the verifier,
and the router before any training, is worth about three and a half points of exact
match (nearly nine on the training benchmarks), as \Cref{fig:headline-decomp} showed.
\end{precise}

\begin{precise}[The no-self-improvement lesion is the cold-start cycle]
One of the three registered lesions needs no separate run, because it already exists in
the trajectory. The no-self-improvement lesion is defined as the full pipeline with
fine-tuning switched off, which is exactly the cold-start cycle of \Cref{ch:trajectory}.
Pricing self-improvement is therefore the cold-start-versus-third-cycle contrast, and
its reading is the one \Cref{ch:benchmarks} and \Cref{ch:trajectory} already established:
the loop is close to flat on pooled exact match and worth about a fifteen percent
reduction in the hallucination metric, a quarter on the training benchmarks. The lesion
confirms what the metric suite was built to surface: self-improvement buys honesty, not
raw accuracy, and removing it costs honesty rather than accuracy.
\end{precise}

\section{The retrieval-utility classifier, switched off}
\label{sec:abl-ruc}

One component was toggled cleanly on and off under a matched protocol: the
retrieval-utility classifier of \Cref{ch:ruc}. It is the chapter's one true
component-level ablation, and it is also the one with the most to teach, because the
result is asymmetric.

\begin{table}[t]
\centering
\small
\begin{tabular}{@{}lccc@{}}
\toprule
\textbf{Benchmark} & \textbf{EM off $\to$ on} & \textbf{CHM off $\to$ on (\% rel)} & \textbf{retrieval share off $\to$ on} \\
\midrule
fact-checking      & $0.503 \to 0.503$ ($+0.0$) & $0.117 \to 0.120$ ($+2.5$)  & $92.3\% \to 70.7\%$ \\
open-text factoid  & $0.437 \to 0.463$ ($+2.7$) & $0.123 \to 0.121$ ($-2.0$)  & $95.7\% \to 68.0\%$ \\
commonsense        & $0.603 \to \mathbf{0.743}$ ($\mathbf{+14.0}$) & $0.131 \to \mathbf{0.086}$ ($\mathbf{-34.4}$) & $57.7\% \to 20.3\%$ \\
implicit multi-step& $0.617 \to 0.640$ ($+2.3$) & $0.119 \to 0.128$ ($+8.1$)  & $100\% \to 76.7\%$ \\
misconception      & $0.210 \to \mathbf{0.310}$ ($\mathbf{+10.0}$) & $0.116 \to 0.119$ ($+2.9$) & $99.0\% \to \mathbf{40.0\%}$ \\
\midrule
\textbf{pooled}    & $\mathbf{0.474 \to 0.532}$ ($\mathbf{+5.8}$) & $\mathbf{0.121 \to 0.115}$ ($\mathbf{-5.2}$) & $\mathbf{88.9\% \to 55.1\%}$ \\
\bottomrule
\end{tabular}
\caption[The retrieval-utility classifier, on versus off]{The actual on-versus-off
contrast for the retrieval-utility classifier at the third cycle, matched backbone and
evaluation fold (misconception exact match by the language-model judge). With the
classifier off, every query that fails the cheap-path gates routes unconditionally to
retrieval; with it on, those queries route to retrieval only when the classifier
predicts retrieval will help. The asymmetric-rescue pattern is in the bold cells.}
\label{tab:abl-ruc}
\end{table}

\begin{figure}[t]
\centering
\begin{tikzpicture}
\begin{axis}[
  width=12.2cm, height=5.4cm, ybar, bar width=8pt,
  symbolic x coords={fact-check, factoid, commonsense, multi-step, misconception},
  xtick=data, x tick label style={font=\scriptsize, rotate=20, anchor=east},
  ymin=0, ymax=0.82, ylabel={exact match}, ytick={0,0.2,0.4,0.6,0.8},
  font=\small\sffamily, tick label style={font=\scriptsize},
  legend style={at={(0.5,1.16)}, anchor=north, legend columns=2, draw=cbInk!25, font=\scriptsize},
  ymajorgrids, grid style={cbInk!8},
]
  \addplot[draw=cbInk!55, fill=cbInk!14] coordinates {
    (fact-check,0.503)(factoid,0.437)(commonsense,0.603)(multi-step,0.617)(misconception,0.210)};
  \addlegendentry{classifier off}
  \addplot[draw=cbExample!70!black, fill=cbExample!28] coordinates {
    (fact-check,0.503)(factoid,0.463)(commonsense,0.743)(multi-step,0.640)(misconception,0.310)};
  \addlegendentry{classifier on}
\end{axis}
\end{tikzpicture}
\caption[The asymmetric rescue]{The actual per-benchmark exact match with the
retrieval-utility classifier off and on. The gains are concentrated, not uniform: the
two benchmarks where unconditional retrieval was hurting (commonsense $+14.0$ points,
misconception $+10.0$ points) are rescued, while the fact-checking benchmark, where
retrieval helps, is held exactly flat. The classifier learns where grounding earns its
cost and routes accordingly, collapsing the pooled retrieval share from $88.9\%$ to
$55.1\%$.}
\label{fig:abl-ruc}
\end{figure}

\begin{precise}[Asymmetric rescue, and why it is the right shape]
The classifier lifts pooled exact match by $5.8$ points and lowers the pooled
hallucination metric by about five percent, but the pooled numbers understate the
mechanism. The gains are concentrated where retrieval was hurting: the commonsense
benchmark gains fourteen points of exact match and loses a third of its hallucination
metric, and the misconception benchmark gains ten points, while the fact-checking
benchmark, where retrieval genuinely helps, is held flat. The retrieval share collapses
from eighty-nine percent to fifty-five percent pooled, and by fifty-nine points on the
misconception benchmark specifically. This is the right shape because it shows the
classifier is not simply turning retrieval down everywhere; it is learning the
per-benchmark sign of retrieval's value and acting on it, preserving grounding where it
pays and withdrawing it where it amplifies the wrong answer.
\end{precise}

\begin{defence}[If an examiner asks: ``is this a clean classifier-only ablation?'']
No, and the chapter says so in the same breath as the result. The on-arm shipped the
classifier alongside concurrent pipeline corrections and a re-fitted composite, so the
contrast is a \emph{bundled} on-versus-off comparison, not a single-variable isolation.
The honest claim is that the bundle, with the classifier as its load-bearing addition,
produces the asymmetric rescue; the arm-isolating re-run that would attribute the gain
to the classifier alone is registered as future work. Reporting the confound next to the
number is the difference between an ablation and a result dressed up as one. The
mechanism, routing retrieval away from the benchmarks it hurts, is exactly what the
classifier was designed to do, which is why the bundled gain is attributed to it even
without the clean isolation.
\end{defence}

\section{The forgetting guard, ablated by the run itself}
\label{sec:abl-guard}

The no-forgetting-guard lesion would set the retention tolerance to zero and let every
update commit. It was not run as a separate sweep, but the run performed the experiment
for us, once, at the fourth cycle.

\begin{precise}[The fourth cycle is the guard's natural ablation]
At the fourth cycle the freshly fine-tuned model reached the lowest training loss of the
entire run, about $0.07$, yet its held-out open-text factoid probe fell to $0.886$,
below the floor. The guard aborted the cycle and rolled the weights back. With the
forgetting guard ablated, that same update would have committed: the model with the
best training loss and a failed retention probe would have become the next cycle's
starting point. The guard is therefore worth exactly the degradation it refused to
accept: an update that looked best by training loss and worst by held-out retention.
This is the cleanest possible illustration of why training loss is not the quantity to
trust, and why the guard reads a held-out probe instead. The clean lesion, which would
let the loop run several cycles without the guard and measure the cumulative
forgetting, would change which cycles commit and so needs its own run; it is registered
but not yet performed. What the run shows is the guard catching a real degrading update
on the one occasion it appeared.
\end{precise}

\section{The lesion still owed, and the honest ledger}
\label{sec:abl-owed}

One registered lesion has neither a direct metric nor a natural observation, and the
chapter does not pretend otherwise.

\begin{pitfall}[No-retroverify was not run]
The no-retroverify lesion skips the per-cycle re-scoring of stored memory. Unlike
no-self-improvement, it cannot be read off an existing cycle, because its effect is a
\emph{divergence} that compounds over cycles: without re-verification the stored pool
would slowly accumulate episodes that a better model would no longer trust, and only a
parallel run under the lesion would measure how far the purity drifts. That run was not
performed. What the trajectory does show is the work re-verification did when it ran: it
pruned a declining fraction of the pool each cycle, about thirteen percent at the first
committed cycle down to about four percent at the last, which is the count of episodes
the re-scoring removed. That is evidence the mechanism is active and that its workload
shrinks as the pool's quality rises, but it is not the counterfactual purity-without-it.
The clean lesion is registered future work, and the book reports it as owed rather than
quietly folding the prune-rate trace into a claim it cannot support.
\end{pitfall}

\begin{table}[t]
\centering
\small
\begin{tabular}{@{}p{3.0cm}p{3.6cm}p{7.2cm}@{}}
\toprule
\textbf{Component} & \textbf{Evidence} & \textbf{What it is worth} \\
\midrule
Three-tier architecture     & baseline panel, cold start            & $+3.4$ pp pooled EM ($+8.8$ training); routing earns the accuracy \\
Self-improvement loop       & cold-start vs trained                 & flat EM, about $-15\%$ CHM (honesty, not accuracy) \\
Retrieval-utility classifier& on/off, bundled                       & $+5.8$ pp EM, $-5.2\%$ CHM, retrieval share $-33.8$ pp \\
Forgetting guard            & cycle-four abort                      & caught a lowest-loss, failed-retention update \\
Retroverification           & prune-rate trace (counterfactual owed)& removed $13\%\!\to\!4\%$ of the pool per cycle \\
\bottomrule
\end{tabular}
\caption[The component ledger]{What each piece is worth, by the cheapest sound evidence
available. Three components are priced by a metric already in hand, one by a clean
on/off toggle (bundled), one by a natural observation, and one by an active-but-not-yet
-counterfactually-isolated trace. The exact-match and hallucination deltas are pooled
over the five-benchmark panel.}
\label{tab:abl-ledger}
\end{table}

\begin{bythenumbers}[The ledger, read honestly]
Pulling the pieces together: the three-tier architecture earns the accuracy
(about three and a half pooled points over the floor, nearly nine on the training
benchmarks); the self-improvement loop earns the honesty (roughly a fifteen percent
reduction in confident wrongness while accuracy stays flat); the retrieval-utility
classifier earns a further five-and-a-half points of accuracy by routing retrieval only
where it helps, though that figure is bundled with concurrent fixes; the forgetting
guard earns its keep on the one cycle it caught a degrading update; and re-verification
is demonstrably active but its clean counterfactual is owed. No single number captures
the system, which is the whole reason the metric suite of \Cref{ch:benchmarks} reports
several. The architecture and the loop divide the labour between them: one makes the
system accurate, the other makes it honest, and the components around them keep both
properties from eroding.
\end{bythenumbers}

This closes the evaluation part. The architecture has been built stage by stage,
proven where proof was possible, measured across a real trajectory, and priced
component by component. \Cref{ch:constants} collects every registered constant the book
has used into one registry, and \Cref{ch:glossary} gathers the symbols and terms, so
that any number or name in the preceding chapters can be traced to its definition.
